\section{Brian Goetz}
Java Language Architect na Oracle; foi spec lead do JSR-335 (lambdas) e é o
autor de \emph{Java Concurrency in Practice} (fonte:
\citeobra{bio:goetz-inside}).

\subsection*{Lambdas e bytecode}
\begin{itemize}
\obra{doc:state-of-lambda}{Lambdas, interfaces funcionais, method
  references e default methods, explicados pelo designer. O documento não
  tem linha de autor: está hospedado no diretório pessoal dele em
  \texttt{cr.openjdk.org/\textasciitilde briangoetz}.}
\obra{doc:lambda-translation}{Como o \texttt{javac} traduz um lambda em
  \texttt{invokedynamic} e \texttt{LambdaMetafactory}. Leia com o
  \texttt{javap -c} aberto. Mesma hospedagem, também sem byline.}
\obra{video:goetz-lambda-hood}{A mesma tradução em palestra.}
\obra{doc:jep126}{O JEP que trouxe os lambdas ao Java~8 (Owner: Goetz).}
\obra{doc:jep484}{A API oficial para ler e gerar bytecode, que substitui o
  ASM dentro do JDK (autor: Goetz).}
\end{itemize}

\subsection*{Java moderno}
\begin{itemize}
\obra{artigo:dop-java}{Records, sealed types e pattern matching como
  modelagem de dados: tornar estados ilegais irrepresentáveis.}
\obra{doc:beyond-records}{Design note do Amber sobre levar a
  desestruturação a classes que não são records.}
\obra{doc:on-ramp}{Por que o Java simplificou o \enquote{hello world}.
  Ótimo para quem está começando.}
\obra{doc:jep286}{O \texttt{var}: motivação, limites e as alternativas
  descartadas (autor: Goetz).}
\obra{artigo:goetz-virtual-threads}{Virtual threads e o modelo
  thread-per-request.}
\obra{video:goetz-fp-vs-oo}{Funcional e orientado a objetos como
  ferramentas complementares.}
\obra{video:goetz-growing-java}{Como e por que a linguagem evolui (JVMLS
  2025).}
\end{itemize}

\subsection*{Valhalla e o futuro da JVM}
\begin{itemize}
\obra{doc:state-of-valhalla}{Value objects, layout de memória e por que
  objetos têm identidade.}
\obra{doc:defense-of-erasure}{Por que os genéricos do Java usam erasure,
  explicado por quem decide o futuro deles.}
\obra{video:goetz-valhalla-refactor}{Value classes, tipos null-restricted
  e strict initialization (Devoxx 2024).}
\end{itemize}
